
\documentclass{article}
\usepackage{colortbl}
\usepackage{makecell}
\usepackage{multirow}
\usepackage{supertabular}

\begin{document}

\newcounter{utterance}

\centering \large Interaction Transcript for game `referencegame', experiment `line\_grids\_rows', episode 11 with reflexion{-}Qwen3.5{-}4B.
\vspace{24pt}

{ \footnotesize  \setcounter{utterance}{1}
\setlength{\tabcolsep}{0pt}
\begin{supertabular}{c@{$\;$}|p{.15\linewidth}@{}p{.15\linewidth}p{.15\linewidth}p{.15\linewidth}p{.15\linewidth}p{.15\linewidth}}
    \# & $\;$A & \multicolumn{4}{c}{Game Master} & $\;\:$B\\
    \hline

    \theutterance \stepcounter{utterance}  
    & & \multicolumn{4}{p{0.6\linewidth}}{
        \cellcolor[rgb]{0.9,0.9,0.9}{
            \makecell[{{p{\linewidth}}}]{
                \texttt{\tiny{[P1$\langle$GM]}}
                \texttt{You are given three grids, where each of them is 5 by 5 in size.} \\
\texttt{Grids have empty cells marked with "▢" and filled cells marked with "X".} \\
\texttt{Your task is to generate a referring expression that best describes the target grid while distinguishing it from the two other distractor grids.} \\
\texttt{The first grid is the target grid, and the following two grids are the distractors.} \\
\\ 
\texttt{Target grid:} \\
\\ 
\texttt{□ □ □ X X} \\
\texttt{□ □ □ X X} \\
\texttt{□ □ □ X X} \\
\texttt{□ □ □ X X} \\
\texttt{□ □ □ X X} \\
\\ 
\texttt{Distractor grid 1:} \\
\\ 
\texttt{□ □ X X X} \\
\texttt{□ □ X X X} \\
\texttt{□ □ X X X} \\
\texttt{□ □ X X X} \\
\texttt{□ □ X X X} \\
\\ 
\texttt{Distractor grid 2:} \\
\\ 
\texttt{X □ □ X X} \\
\texttt{X □ □ X X} \\
\texttt{X □ □ X X} \\
\texttt{X □ □ X X} \\
\texttt{X □ □ X X} \\
\\ 
\texttt{Instruction: Describe the target grid.} \\
\texttt{Generate the referring expression starting with the tag "Expression: " for the given target grid. Omit any other text.} \\
            }
        }
    }
    & \\ \\

    \theutterance \stepcounter{utterance}  
    & \multicolumn{4}{p{0.6\linewidth}}{
        \cellcolor[rgb]{0.9,0.9,0.9}{
            \makecell[{{p{\linewidth}}}]{
                \texttt{\tiny{[P1$\rangle$GM]}}
                \texttt{Expression: the grid where the filled cells form a 2x5 block in the top{-}right corner} \\
            }
        }
    }
    & & \\ \\

    \theutterance \stepcounter{utterance}  
    & & & \multicolumn{2}{p{0.3\linewidth}}{
        \cellcolor[rgb]{0.9,0.9,0.9}{
            \makecell[{{p{\linewidth}}}]{
                \texttt{\tiny{[GM$|$GM]}}
                \texttt{Expression: the grid where the filled cells form a 2x5 block in the top{-}right corner} \\
            }
        }
    }
    & & \\ \\

    \theutterance \stepcounter{utterance}  
    & & \multicolumn{4}{p{0.6\linewidth}}{
        \cellcolor[rgb]{0.9,0.9,0.9}{
            \makecell[{{p{\linewidth}}}]{
                \texttt{\tiny{[GM$\rangle$P2]}}
                \texttt{You are given three grids, where each of them is 5 by 5 in size.} \\
\texttt{Grids have empty cells marked with "▢" and filled cells marked with "X".} \\
\texttt{You are also given a referring expression that describes one of the given grids.} \\
\texttt{Your task is to select the grid that matches the given referring expression.} \\
\texttt{Generate only the number (in text) of the grid that matches the given expression by selecting first, second, or third.} \\
\\ 
\texttt{First:} \\
\\ 
\texttt{X □ □ X X} \\
\texttt{X □ □ X X} \\
\texttt{X □ □ X X} \\
\texttt{X □ □ X X} \\
\texttt{X □ □ X X} \\
\\ 
\texttt{Second:} \\
\\ 
\texttt{□ □ X X X} \\
\texttt{□ □ X X X} \\
\texttt{□ □ X X X} \\
\texttt{□ □ X X X} \\
\texttt{□ □ X X X} \\
\\ 
\texttt{Third:} \\
\\ 
\texttt{□ □ □ X X} \\
\texttt{□ □ □ X X} \\
\texttt{□ □ □ X X} \\
\texttt{□ □ □ X X} \\
\texttt{□ □ □ X X} \\
\\ 
\texttt{Expression: the grid where the filled cells form a 2x5 block in the top{-}right corner} \\
\texttt{Question: Which grid does the expression refer to?} \\
\texttt{Start with the tag "Answer: ", followed by your selection. Omit any other text.} \\
            }
        }
    }
    & \\ \\

    \theutterance \stepcounter{utterance}  
    & & & \multicolumn{4}{p{0.6\linewidth}}{
        \cellcolor[rgb]{0.9,0.9,0.9}{
            \makecell[{{p{\linewidth}}}]{
                \texttt{\tiny{[GM$\langle$P2]}}
                \texttt{Answer: First} \\
            }
        }
    }
     \\ \\

    \theutterance \stepcounter{utterance}  
    & & & \multicolumn{2}{p{0.3\linewidth}}{
        \cellcolor[rgb]{0.9,0.9,0.9}{
            \makecell[{{p{\linewidth}}}]{
                \texttt{\tiny{[GM$|$GM]}}
                \texttt{first} \\
            }
        }
    }
    & & \\ \\

    \theutterance \stepcounter{utterance}  
    & \multicolumn{4}{p{0.6\linewidth}}{
        \cellcolor[rgb]{0.95,0.95,0.95}{
            \makecell[{{p{\linewidth}}}]{
                \texttt{\tiny{[P2$\rightarrow$P1]}}
                \texttt{Post{-}episode Reflexion (not sent to GM):} \\
\texttt{New plan: The previous expression "2x5 block" was ambiguous because both the target grid (right{-}aligned) and distractor grid 1 (left{-}aligned) consist of a 2x5 block of filled cells. To distinguish the target grid, I must specify the horizontal position of the block relative to the grid boundaries. Next time, I will explicitly state that the filled cells are in the "rightmost" or "right{-}hand" columns, or describe the empty space as being in the "leftmost" columns, ensuring the description uniquely identifies the target among the distractors.} \\
            }
        }
    }
    & & \\ \\

    \theutterance \stepcounter{utterance}  
    & & & \multicolumn{4}{p{0.6\linewidth}}{
        \cellcolor[rgb]{0.95,0.95,0.95}{
            \makecell[{{p{\linewidth}}}]{
                \texttt{\tiny{[P1$\rightarrow$P2]}}
                \texttt{Post{-}episode Reflexion (not sent to GM):} \\
\texttt{New plan: I misidentified the target grid by focusing on the wrong visual feature. The expression specifies a "2x5 block in the top{-}right corner," which means the filled cells should occupy the top two rows and the rightmost five columns. I incorrectly selected the First grid, which has filled cells in the leftmost column. I should have analyzed the position of the filled "X" blocks relative to the grid edges: the correct grid (Second) has filled cells in the top{-}right, while the First has them on the left. Next time, I will strictly verify that the filled cells form a rectangle of height 2 and width 5 located specifically in the top{-}right quadrant before answering.} \\
            }
        }
    }
     \\ \\

\end{supertabular}
}

\end{document}
